\chapter{相关硬件基础}

\section{中央处理器}
\concept{CPU 的基本组成与工作流程}
\begin{points}
  \item CPU 的主要部件包括控制器、运算器、寄存器和高速缓存。
  \item CPU 的任务是执行指令，基本流程为：取指 $\rightarrow$ 译码 $\rightarrow$ 取操作数 $\rightarrow$ 执行指令。
  \item 影响 CPU 性能的因素包括主频、字长、指令集、高速缓存、核心数量、总线速度和处理技术。
\end{points}
\plain{这里先抓 CPU 怎样一条条执行指令，后面理解中断、系统调用、上下文切换都要靠这条执行主线。}
\exam{常考取指、译码、执行的顺序，以及 CPU 性能相关因素的辨析。}


\concept{单处理器与多处理器}
\begin{points}
  \item 单处理器系统只有一个运算处理器，同一时刻只能有一个进程在 CPU 上运行。
  \item 单 CPU 中，进程与进程只能并发不能并行；处理机与设备、设备与设备在硬件支持下可以并行工作。
  \item 多处理器系统有多个运算处理器，可分为对称多处理 SMP 和非对称多处理 AMP。
  \item SMP 中各处理器地位相同，共享内存和 I/O 系统；AMP 中主从处理器角色不同，主处理器负责调度与管理。
\end{points}
\plain{单处理器重点分清“并发但不并行”，多处理器才可能多个进程真的同时在不同 CPU 上跑。}
\exam{常考单处理机和多处理机各能实现什么并行关系，以及 SMP、AMP 的区别。}


\concept{Flynn 分类}
\begin{points}
  \item SISD：单指令流单数据流，传统顺序单处理器。
  \item SIMD：单指令流多数据流，一条指令同时处理多个数据，如向量处理、GPU 中的某些模式。
  \item MISD：多指令流单数据流，实际很少见。
  \item MIMD：多指令流多数据流，多处理器/多核系统常见。
\end{points}
\plain{Flynn 分类本质是在看“几条指令流、几条数据流”同时工作，不要和单核多任务混在一起。}
\exam{常考 SISD、SIMD、MIMD 的定义和典型代表。}


\concept{多处理器、多核与超线程}
\begin{points}
  \item 多处理器通常指在一个系统中放置多个 CPU 芯片；多核通常指在同一块 CPU 芯片内集成多个执行核。
  \item 多核结构通常比多 CPU 更紧凑，成本更低、功耗更小。
  \item 超线程是在一个物理处理器内部提供多个逻辑处理器/线程执行环境，共享执行部件和缓存，让处理器更充分利用流水线和功能部件。
  \item 直觉：多核更像“真的多套核心”；超线程更像“一套核心尽量同时喂多条线程”。
\end{points}
\plain{题目若把多核和超线程放一起，抓“多核是真多执行核，超线程是共享执行部件的多个逻辑执行环境”就够了。}
\exam{常考多处理器、多核、超线程的区别，以及谁是真正增加了完整执行核心。}


\concept{操作系统与多核处理器}
\begin{points}
  \item 多核处理器要求操作系统在处理器通信、数据共享、存储层次管理、同步、调度优化和能耗管理等方面提供支持。
  \item 核数增加后，系统不只是“多几个 CPU”，还要处理缓存一致性、线程同步、负载均衡等问题。
  \item 因此 OS 需要改进调度策略、并行执行模型和共享数据管理机制。
\end{points}
\plain{多核不是硬件自己多加几个核就结束了，OS 还得负责把线程分好、数据管好、同步做好。}
\exam{常考操作系统为何要针对多核做调度、同步和共享支持。}


\section{寄存器、特权指令与处理器状态}
\concept{寄存器}
\begin{points}
  \item 寄存器是 CPU 内部高速存储单元，用于保存指令、地址、数据和状态信息。
  \item 常见寄存器包括程序计数器 PC、指令寄存器 IR、通用寄存器、栈指针 SP、程序状态字 PSW。
  \item 上下文切换时，操作系统需要保存当前进程的寄存器现场，并恢复下一个进程的现场。
\end{points}
\plain{寄存器是最直接的执行现场；只要发生中断或进程切换，OS 首先得保存和恢复这些内容。}
\exam{常考 PC、IR、SP、PSW 等寄存器的作用，以及哪些内容属于处理机现场。}


\concept{寄存器的功能分类}
\begin{points}
  \item 以 x86 为例，寄存器可分为通用寄存器、指针及变址寄存器、段寄存器、指令指针与标志寄存器、控制寄存器等。
  \item 通用寄存器常用于普通算术逻辑和数据暂存；栈指针、基址寄存器等常用于过程调用和栈管理。
  \item 段寄存器、控制寄存器、标志寄存器更偏向地址变换、保护和处理器控制。
  \item 外设控制器也可能带有设备专用寄存器，但它们不等同于 CPU 内部通用寄存器。
\end{points}
\plain{这类题通常不是让你死背所有寄存器名，而是分清哪些寄存器管普通数据，哪些管控制、状态、地址。}
\exam{常考按功能给寄存器分类，或判断某寄存器更偏现场信息还是控制信息。}


\concept{特权指令与非特权指令}
\begin{points}
  \item 非特权指令可由用户程序直接执行，如普通算术、逻辑、数据移动指令。
  \item 特权指令只能由内核执行，如 I/O 指令、关中断、修改内存保护寄存器、修改页表、启动 DMA、设置定时器。
  \item 判断标准：凡是可能影响系统安全、影响其他进程、直接控制硬件资源的操作，通常属于特权操作。
\end{points}
\plain{普通程序不能直接碰全局危险资源；一旦需要 I/O、页表、中断等操作，就要进内核态。}
\exam{常考哪些指令是特权指令、系统调用是否引起模式切换、模式切换和进程切换区别。}


\concept{特权指令典型操作}
\begin{points}
  \item 典型特权操作包括设置定时器、读取时钟、清除内存、发起陷入、关中断、修改设备状态、切换用户态与内核态、访问 I/O 设备等。
  \item 这类操作一旦被普通用户程序任意执行，就可能破坏调度、公平性、保护机制和设备控制。
  \item 所以题目判断时，只要操作能直接改变系统控制权或全局硬件状态，通常就应归入特权操作。
\end{points}
\plain{判断特权指令时别死记列表，直接想：这事如果让普通程序自己干，会不会把整个系统搞乱。}
\exam{常考举例判断某操作是否属于特权操作，并说明原因。}


\concept{处理器状态：用户态与内核态}
\begin{points}
  \item 用户态下执行普通用户程序，只能执行非特权指令，不能直接访问受保护资源。
  \item 内核态下执行操作系统内核，可执行特权指令，访问核心数据结构。
  \item 从用户态到内核态的常见原因：系统调用、异常、外部中断。
  \item 从内核态返回用户态：中断/系统调用处理完成后，恢复用户进程现场并执行返回指令。
\end{points}
\plain{用户态和内核态本质是硬件给 OS 划出的权限边界，危险操作必须越过这条边界让内核代办。}
\exam{常考哪些事件会从用户态进入内核态，以及模式切换是否一定伴随进程切换。}


\concept{保护级别与 Ring}
\begin{points}
  \item 一些处理器不只区分两态，还可支持多个保护级别。
  \item 例如 Intel Pentium 体系结构有 4 个保护级别，Ring 0 权限最高，Ring 3 权限最低。
  \item 内核通常运行在 Ring 0，普通用户程序通常运行在 Ring 3。
  \item 这本质上仍是在做权限隔离，只是比“用户态/内核态”描述得更细。
\end{points}
\plain{考试里看到 Ring，直接把它理解成“更细粒度的权限圈层”；Ring0 最强，Ring3 最弱。}
\exam{常考 Ring0 和 Ring3 分别对应什么，以及它和用户态/内核态的关系。}


\concept{处理器状态切换步骤}
\begin{points}
  \item 用户态进入内核态的原因主要有两类：程序主动发起系统调用，或程序运行中发生中断/异常。
  \item 常见切换步骤包括：保存现场、根据中断号或陷入入口设置程序计数器、交换或修改 PSW、转入对应处理程序。
  \item 处理完成后，再恢复现场并返回原程序或转入调度程序。
\end{points}
\plain{把状态切换看成“先保现场，再找入口，再换权限去处理”，答题就不容易漏步骤。}
\exam{常考用户态切到内核态的触发原因，以及状态切换的基本顺序。}


\concept{PSW 程序状态字}
\begin{points}
  \item PSW 保存 CPU 的运行状态信息，包括条件码、中断屏蔽位、处理器模式位等。
  \item 模式位用于区分用户态和内核态，是硬件支持操作系统保护机制的重要基础。
  \item 中断发生时，硬件通常会保存 PC/PSW 等关键信息，以便中断处理完成后回到原程序继续执行。
\end{points}
\plain{PSW 可以看成 CPU 当前“运行身份和状态”的摘要，模式位、中断位都在这里体现。}
\exam{常考 PSW 的作用，以及它为什么是中断返回和现场恢复的重要内容。}


\concept{x86 总体架构中的关键寄存器}
\begin{points}
  \item 指令指针寄存器 EIP 保存下一条待执行指令的地址，常被 \code{CALL}、\code{RET}、\code{JMP} 等控制转移指令修改。
  \item 控制寄存器中常见位有保护模式位和分页相关位，它们决定处理器是否启用保护和分页机制。
  \item 与内存管理相关的寄存器还包括 GDTR、IDTR、TR 等，分别和全局描述符、中断描述符、任务状态保存等机制相关。
  \item 这些细节本章只需知道“它们服务于保护、分页、中断和任务切换”，不必按体系结构课深度去背位级细节。
\end{points}
\plain{x86 这部分不是要你变成体系结构专家，而是知道哪些寄存器在支撑 OS 的保护、中断和任务切换。}
\exam{常考 EIP、控制寄存器、GDTR/IDTR/TR 大概各管什么。}


\section{中断处理}
\concept{中断的基本概念}
\begin{points}
  \item 中断是 CPU 暂停当前程序执行，转而执行相应处理程序，处理完成后再返回原程序的机制。
  \item 中断使 CPU 不必忙等 I/O 完成，是实现多道程序设计和 CPU/I/O 并行工作的关键。
  \item 中断来源包括外部设备中断、定时器中断、异常和系统调用陷入。
\end{points}
\plain{答题时把步骤串起来：提出请求、保存现场、查中断向量、执行服务程序、恢复现场。}
\exam{常考中断响应流程排序题，或简答中断处理的基本步骤。}

\concept{中断分类}
\begin{points}
  \item \term{异步中断}通常来自 CPU 外部，如定时器中断、I/O 设备中断，与当前正在执行的指令本身无直接因果关系。
  \item \term{同步中断}通常来自 CPU 内部，也常称异常，由当前指令执行直接引起。
  \item 在不少体系结构表述中，系统调用也通过一种特殊异常/陷入机制实现。
  \item 因而教材里有时把“中断”狭义地只指异步外部中断，有时又把中断与异常统称为广义中断。
\end{points}
\plain{这块最容易混的是“中断”这个词有狭义和广义两层用法；做题时先看它是在按来源分，还是在把异常一起总称。}
\exam{常考定时器/I/O 属于异步中断，除零/缺页属于同步异常，以及系统调用在实现上通常归入哪一类。}


\concept{中断处理过程}
\begin{points}
  \item 设备或内部事件提出中断请求。
  \item CPU 在适当时刻响应中断，保存断点和现场，切换到内核态。
  \item 根据中断向量找到对应的中断服务程序。
  \item 中断服务程序处理中断事件，例如读取设备状态、唤醒等待进程。
  \item 恢复被中断程序现场，返回原程序或转入调度程序。
\end{points}
\plain{中断就是硬件或软件打断 CPU 当前执行，让 OS 先处理更紧急或必须处理的事件。}
\exam{常考中断处理步骤、中断与异常/系统调用区别、为什么中断能提高 CPU 与 I/O 并行度。}

\concept{中断装置与中断处理程序}
\begin{points}
  \item 发现中断源、仲裁优先级并引导 CPU 转入相应服务程序的硬件机制，可统称为\term{中断装置}。
  \item 当多个中断同时请求时，中断装置要先识别中断源并按优先级选择先响应者。
  \item \term{中断处理程序}负责具体处理中断事件，如读取设备状态、完成数据收尾、唤醒等待进程，并在结束后恢复现场。
  \item 因而可以把中断装置理解为“把 CPU 正确带到入口”，把中断处理程序理解为“真正把事情办完”。
\end{points}
\plain{考试里不要把硬件发现中断源和软件处理中断事件混成一步：前者偏“找到谁在叫、跳到哪去”，后者偏“具体怎么收尾”。}
\exam{若简答题问中断装置或中断处理程序的作用，答题时最好把“识别中断源/优先级选择/转入入口”和“处理事件/恢复现场”分开写。}


\concept{中断与异常、系统调用}
\begin{points}
  \item 外部中断通常来自 CPU 外部设备，如磁盘、网卡、键盘。
  \item 异常来自当前指令执行过程，如除零、非法指令、缺页异常。
  \item 系统调用是用户程序有意执行陷入指令，请求内核服务。
  \item 三者都会导致从用户态进入内核态，但触发原因不同。
\end{points}
\plain{三者都会陷入内核，但来源不同：外部设备来的是中断，当前指令出错是异常，程序主动请求服务是系统调用。}
\exam{常考中断、异常、系统调用的来源、同步异步性和是否由程序主动发起。}


\concept{中断优先级与多重中断}
\begin{points}
  \item 当多个中断请求同时到来时，系统通常按预先规定的优先级决定先响应谁。
  \item 高优先级中断可在一定条件下打断低优先级中断服务程序，这就形成多重中断。
  \item 是否允许中断嵌套，与处理中断时是否重新开中断、硬件优先级控制和现场保护能力有关。
  \item 多重中断提高了系统对紧急事件的响应能力，但也增加了中断处理和现场保存恢复的复杂度。
\end{points}
\plain{多重中断可以理解成“处理中断时又被更紧急的中断插队”，前提是系统能分清优先级并正确保存多层现场。}
\exam{PPT 单独列了这节，常考为什么需要中断优先级、什么情况下会出现中断嵌套，以及它带来的代价。}


\concept{同步与原子性}
\begin{points}
  \item \term{同步}强调多个事件或操作之间存在先后配合关系，必须按某种时序推进。
  \item \term{原子性}强调一个操作执行期间不可被打断，要么全部完成，要么完全不做。
  \item 在操作系统里，关中断、硬件原语、原子读改写指令等常用于保证关键操作的原子性。
  \item 原子性是后续实现互斥和同步机制的重要硬件基础，但“同步”和“原子”不是同一个概念。
\end{points}
\plain{同步管“先后顺序要配合”，原子性管“中间不能被拆开”；两个词经常一起出现，但问的不是一回事。}
\exam{这节虽然在硬件基础里，但会直接服务后面进程同步章节，常考同步与原子性的区别以及为什么原子指令对互斥实现重要。}


\section{操作系统的硬件支持}
\concept{内存保护}
\begin{points}
  \item 为防止进程访问非法内存，硬件需要支持地址变换和地址合法性检查。
  \item 简单系统中可用基址寄存器和界限寄存器检查地址是否越界。
  \item 分页/分段系统中由 MMU 根据页表或段表完成地址映射和保护检查。
\end{points}
\plain{内存保护的核心是“一个进程不能随便碰另一个进程或内核的地址空间”，否则多道程序根本不安全。}
\exam{常考基址/界限寄存器或页表、段表如何支持地址检查与保护。}


\concept{定时器}
\begin{points}
  \item 定时器用于防止用户程序长期独占 CPU。
  \item 操作系统设置时间片，到时后定时器中断发生，CPU 转入内核执行调度程序。
  \item 分时系统和抢占式调度离不开定时器支持。
\end{points}
\plain{定时器是抢占式 OS 的底盘，没有时钟中断，系统就很难防止用户程序长期霸占 CPU。}
\exam{常考定时器为什么是分时系统和抢占式调度的硬件基础。}


\concept{I/O 保护}
\begin{points}
  \item 用户程序不能直接执行 I/O 指令，必须通过系统调用请求内核完成。
  \item 这样可以统一设备分配、访问权限、缓冲、错误处理，防止多个进程直接抢设备。
  \item 因此 I/O 指令属于典型特权指令。
\end{points}
\plain{I/O 保护强调的是“普通进程不能直接发设备命令”，否则设备分配、权限和安全性都会失控。}
\exam{常考为什么 I/O 指令是特权指令，以及用户程序怎样合法发起 I/O。}


\section{计算机启动与虚拟机}
\concept{x86 实模式与保护模式}
\begin{points}
  \item 8086 是 16 位处理器，但地址总线为 20 位，因此通过“段基址 + 偏移”完成寻址，实模式地址可写为 $\text{physical addr}=16\times \text{segment}+\text{offset}$。
  \item 80386 及后续处理器为兼容早期体系结构，通常先从实模式启动，再切换到保护模式。
  \item 保护模式支持更强的地址保护、权限管理和更大的寻址能力，是现代操作系统正常运行的基础。
  \item 直觉：实模式更像早期裸奔方式，保护模式才真正支持 OS 要的保护与管理。
\end{points}
\plain{实模式重点看“段基址+偏移”的老式寻址；保护模式重点看“支持现代 OS 的保护机制”。}
\exam{常考为什么 x86 要先实模式启动再切到保护模式，以及实模式的基本寻址公式。}


\concept{计算机启动过程}
\begin{points}
  \item 上电后 CPU 从固定地址开始执行固件中的引导程序。
  \item 固件进行硬件自检和基本初始化，找到可启动设备。
  \item 引导程序把操作系统内核装入内存。
  \item 内核初始化数据结构、设备驱动、内存管理、进程管理，创建第一个用户态进程或系统服务。
\end{points}
\plain{开机不是直接进 OS，而是固件先跑，找到引导程序，再把 OS 内核装入内存。}
\exam{常考 BIOS/UEFI、引导程序、内核加载的大致顺序。}


\concept{虚拟机}
\begin{points}
  \item 虚拟机把底层物理机器抽象成多台逻辑机器，使多个 OS 能在同一物理平台上运行。
  \item 虚拟机监控器负责 CPU、内存、设备等资源的虚拟化与隔离。
  \item 直觉：普通 OS 把裸机虚拟成“适合应用使用的机器”；虚拟机监控器把裸机虚拟成“多台可运行 OS 的机器”。
\end{points}
\plain{虚拟机和“OS 的虚拟资源”不是一回事，这里是再加一层软件把一台物理机伪装成多台可装 OS 的机器。}
\exam{常考虚拟机监控器和普通操作系统的区别，以及虚拟机的基本作用。}
